\documentclass{article}

\input{header}

\title{CSCI 2590: Assignment 4}
\author{Yvonne Yu \\ yy5919 \\ \href{https://github.com/Yvonne-Yu217/CS2590_hw4_yy_1}{https://github.com/Yvonne-Yu217/CS2590\_hw4\_yy\_1}}
\date{}

\colmfinalcopy
\begin{document}
\maketitle

\section*{Question Map}
To make grading easier, this report is organized as:
\begin{itemize}
  \item \textbf{Part I}: Q1, Q2, Q3
  \item \textbf{Part II}: Q4, Q5, Q6, Q7
  \item \textbf{Submission requirements}: repository link, files checklist, and AI usage
\end{itemize}

\section*{Repository Link (Submission Requirement)}
GitHub repository for both Part I and Part II:
\href{https://github.com/Yvonne-Yu217/CS2590_hw4_yy_1}{https://github.com/Yvonne-Yu217/CS2590\_hw4\_yy\_1}

\section*{Q1 Fine-tuning BERT Model}
I implemented the missing training loop in \texttt{do\_train} (forward, loss backpropagation,
optimizer step, scheduler step, gradient reset, and epoch loss reporting), then trained BERT on IMDB.
On the full training/evaluation run, the original test accuracy is \textbf{0.9237}, and the required
output file \texttt{out\_original.txt} is generated.

\section*{Q2 Data Transformations}
I use a token-level stochastic transformation that combines synonym replacement with typo injection.
For each alphabetic token, with probability $p=0.35$, I attempt to transform it. If WordNet provides
valid synonyms different from the original form, I replace the token with a randomly chosen synonym
(with 70\% probability); otherwise (or with 30\% probability even when synonyms exist), I inject a typo
by randomly choosing one of three operations: adjacent character swap, keyboard-neighbor substitution,
or character deletion (for tokens longer than 3 characters).
Non-alphabetic tokens are left unchanged.

This is a reasonable transformation because it simulates realistic user variation: paraphrasing and
small typing mistakes. In most cases sentiment should be preserved, so it serves as an OOD robustness test
without changing labels.

\section*{Q3 Data Augmentation}
\textbf{Accuracy results}
\begin{itemize}
  \item Base model on original test: 0.9237
  \item Base model on transformed test: 0.8367
  \item Augmented model on original test: 0.9287
  \item Augmented model on transformed test: 0.8950
\end{itemize}

\textbf{Analysis}
\begin{itemize}
  \item Performance on transformed data improved significantly after augmentation: $0.8950 - 0.8367 = +0.0583$.
  \item Performance on original data slightly improved: $0.9287 - 0.9237 = +0.0050$.
  \item Intuitively, augmentation teaches the model invariance to surface perturbations while preserving
  core sentiment cues.
  \item One limitation: this augmentation is lexical/local and does not address broader distribution shifts
  such as discourse changes, long-range syntax changes, or sarcasm.
\end{itemize}

\section*{Q4 Data Statistics and Processing}
\begin{table}[h!]
\centering
\begin{tabular}{lcc}
\toprule
Statistics Name & Train & Dev \\
\midrule
Number of examples & 4225 & 466 \\
Mean sentence length (T5 tokens) & 18.10 & 18.07 \\
Mean SQL query length (T5 tokens) & 216.37 & 210.05 \\
Vocabulary size (natural language) & 792 & 466 \\
Vocabulary size (SQL) & 555 & 395 \\
\bottomrule
\end{tabular}
\caption{Data statistics before pre-processing.}
\label{tab:data_stats_before}
\end{table}

\begin{table}[h!]
\centering
\begin{tabular}{lcc}
\toprule
\textbf{T5 fine-tuned model} & Train & Dev \\
\midrule
Model name & google-t5/t5-small & google-t5/t5-small \\
Mean sentence length after prompt (T5 tokens) & 23.10 & 23.07 \\
Mean SQL query length (T5 tokens) & 216.37 & 210.05 \\
Vocabulary size (natural language, prompted) & 796 & 470 \\
Vocabulary size (SQL) & 555 & 395 \\
\bottomrule
\end{tabular}
\caption{Data statistics after pre-processing.}
\label{tab:data_stats_after}
\end{table}

\newpage

\section*{Q5 T5 Fine-tuning Details}
\begin{table}[h!]
\centering
\begin{tabular}{p{3.5cm}p{10cm}}
\toprule
Design choice & Description \\
\midrule
Data processing & Input prompt format: ``translate English to SQL: <NL>''.
Decoder targets are ground-truth SQL sequences plus EOS token. \\
Tokenization & Default \texttt{google-t5/t5-small} tokenizer for both encoder and decoder.
Dynamic padding in collate functions. \\
Architecture & Fine-tuned pretrained T5-small end-to-end (all parameters trainable). \\
Hyperparameters & AdamW, learning rate $3\times10^{-4}$, weight decay 0.01, cosine scheduler, warmup epochs = 2,
batch size = 16, test batch size = 16, max epochs = 30, patience = 5. Generation used beam search
(num\_beams=4, max\_length=512, early\_stopping=True). \\
\bottomrule
\end{tabular}
\caption{Best-performing fine-tuned T5 configuration.}
\label{tab:t5_results_ft}
\end{table}

\section*{Q6 Results and Error Analysis}
\paragraph{Quantitative Results}
\begin{table}[h!]
\centering
\begin{tabular}{lcc}
  \toprule
  System & Query EM & F1 Score \\
  \midrule
  \multicolumn{3}{l}{\textbf{Dev Results}} \\
  \midrule
  T5 fine-tuned (best epoch) & 1.93 & 73.81 \\
  \bottomrule
\end{tabular}
\caption{Development and test results.}
\label{tab:results}
\end{table}

\paragraph{Qualitative Error Analysis}
The following error analysis is for the fine-tuned T5 model (the only model trained, as no extra-credit from-scratch model was attempted).
\begin{table}[h!]
  \centering
  \begin{tabular}{p{2cm}p{3cm}p{5cm}p{3cm}}
    \toprule
    \textbf{Error Type} & \textbf{Example} & \textbf{Description} & \textbf{Statistics} \\
    \midrule
    Syntax error & malformed SQL query & Parentheses, joins, or clause ordering errors causing SQL execution failure. & 56/466 (12.02\%) \\
    \midrule
    Timeout & complex/looping query & Generated query takes too long to execute, often due to redundant joins or Cartesian products. & 8/466 (1.72\%) \\
    \midrule
    Wrong grounding & wrong city/airline/date values & Query is executable but entities or constants differ from the intent, yielding incorrect records. & 83/466 (17.81\%) \\
    \bottomrule
  \end{tabular}
  \caption{Representative error categories on the dev set.}
  \label{tab:qualitative}
\end{table}

\section*{Q7 Checkpoint Link}
Google Drive link containing the model checkpoint used to generate Q7 outputs:
\href{https://drive.google.com/drive/folders/1SdbFj9EizBOGY4g_wG2hFSvgeNt87vIM?usp=sharing}{\url{https://drive.google.com/drive/folders/1SdbFj9EizBOGY4g_wG2hFSvgeNt87vIM?usp=sharing}}

\section*{Submission Files Checklist}
\textbf{Part I (upload to Gradescope)}
\begin{itemize}
  \item out\_original.txt
  \item out\_transformed.txt
  \item out\_augmented\_original.txt
  \item out\_augmented\_transformed.txt
\end{itemize}

\textbf{Part II (upload to Gradescope)}
\begin{itemize}                                                                         
  \item results/t5\_ft\_experiment\_test.sql
  \item records/t5\_ft\_experiment\_test.pkl
\end{itemize}

\textbf{Report (upload to Gradescope)}
\begin{itemize}
  \item hw4-report.pdf
\end{itemize}

\textbf{Google Drive (do not upload to Gradescope)}
\begin{itemize}
  \item checkpoint archive used for Q7 outputs (e.g., v8\_final\_best\_checkpoint.tar.gz)
\end{itemize}

\section*{AI Usage}
\noindent\textbf{1. AI Tools Used}\\
ChatGPT/Copilot-style coding assistant.

\noindent\textbf{2. How AI Was Used}\\
Used for debugging environment/runtime issues, improving training/evaluation script reliability,
and polishing report language. All final code/results were validated by running experiments.

\end{document}
